%% kapa: exam=1 round=23 year=2012 period=1 form=A subject=civil_law subject_ko=민법 seq=27 number=27
%% source: Q-Net 감정평가사 기출문제 (cst00309) artlSeq=5052578 file=raw/1008/5052578/2012년 제23회 감정평가사 제1차시험 기출문제(1교시 A형).pdf
%% conversion: kordoc 4.24.51 -> tools/gen1.py (auto); review=proofread-sonnet; flags=none
\begin{question}{27}
자주점유에 관한 설명으로 옳지 않은 것은? (다툼이 있으면 판례에 의함)
\begin{choices}
  \item 매도인에게 처분권이 없었다는 이유로 매매가 무효가 된 경우 그 사실을 알지 못한 매수인의 점유는 자주점유이다.
  \item 공유자 1인이 공유토지 전부를 점유하고 있는 경우, 특별한 사정이 없는 한 다른 공유자의 지분비율 범위에 대해서는 타주점유에 해당한다.
  \item 등기명의가 신탁되었다면 특별한 사정이 없는 한 명의수탁자의 수탁부동산에 관한 점유는 타주점유이다.
  \item 부동산을 타인에게 매도하여 소유권이전등기를 경료한 후 인도의무를 지고 있는 매도인의 점유는 특별한 사정이 없는 한 타주점유로 변경된다.
  \item 매매대상 토지의 면적이 등기부상의 면적을 상당히 초과하더라도 특별한 사정이 없는 한 초과부분에 대한 매수인의 점유는 자주점유이다.
\end{choices}
\end{question}
\answer{5}
